\section{Linux Investigation Commands}

Linux systems are common targets in enterprise environments, and understanding how to investigate them is essential for CTI analysts, threat hunters, and detection engineers. While SIEM telemetry provides high‑level visibility, direct Linux commands offer immediate insight into processes, users, network activity, and system integrity. This section provides a focused set of commands that support real-world investigations, along with explanations of when and why each command is useful.

Investigations often begin with understanding what is running on a system. The \texttt{ps} and \texttt{top} utilities reveal active processes, their parent-child relationships, and resource usage. These commands help identify suspicious binaries, unexpected services, or processes running under unusual accounts. For example, the following command lists all running processes with full details:

\begin{lstlisting}[language=bash]
ps aux
\end{lstlisting}

To investigate a specific process, analysts often need to inspect its parent process, command-line arguments, and open files. The \texttt{lsof} command is particularly valuable for identifying network connections or files associated with a suspicious process:

\begin{lstlisting}[language=bash]
lsof -p <PID>
\end{lstlisting}

Network visibility is equally important. Commands such as \texttt{ss} and \texttt{netstat} reveal active connections, listening ports, and associated processes. These commands help identify unauthorized services, reverse shells, or unexpected outbound traffic:

\begin{lstlisting}[language=bash]
ss -tulnp
\end{lstlisting}

Authentication and privilege misuse are common indicators of compromise. Linux logs authentication events in \texttt{/var/log/auth.log}, which can be inspected directly or forwarded to Splunk or Sentinel. The following command shows recent authentication attempts:

\begin{lstlisting}[language=bash]
grep "session" /var/log/auth.log
\end{lstlisting}

File integrity is another critical area. Attackers often modify system binaries, create persistence mechanisms, or drop malicious files. The \texttt{find} command helps locate recently modified files, which can reveal tampering or unauthorized activity:

\begin{lstlisting}[language=bash]
find / -type f -mtime -1 2>/dev/null
\end{lstlisting}

System services and startup mechanisms are frequent targets for persistence. The \texttt{systemctl} command provides insight into active services, failed services, and recently modified units:

\begin{lstlisting}[language=bash]
systemctl list-units --type=service
\end{lstlisting}

For deeper forensic analysis, \texttt{auditd} provides detailed system call logging, including file access, privilege escalation attempts, and execution events. When enabled, audit logs can be queried directly:

\begin{lstlisting}[language=bash]
ausearch -m EXECVE
\end{lstlisting}

Linux investigations often require correlating endpoint findings with SIEM telemetry. For example, a suspicious process discovered with \texttt{ps aux} can be cross-referenced with Sysmon for Linux or audit logs ingested into Splunk or Sentinel. The following SPL and KQL examples illustrate how to pivot from endpoint commands to SIEM queries:

\textbf{SPL Example}
\begin{lstlisting}
index=linux_logs process_name="bash"
| table _time, host, user, parent_process, command
\end{lstlisting}

\textbf{KQL Equivalent}
\begin{lstlisting}
LinuxSysmon
| where ProcessName == "bash"
| project TimeGenerated, HostName, User, ParentProcessName, CommandLine
\end{lstlisting}

These combined techniques—endpoint commands plus SIEM queries—form the foundation of Linux threat hunting and incident response. As the handbook progresses, these commands will be integrated into full hunting playbooks and detection patterns tailored to Linux environments.
